% --------------------------------------------------
% 5. RESULTS
% --------------------------------------------------
\section{Results}
\label{sec:results}

\subsection{Basic Cox Model}

Table~\ref{tab:cox_basic} reports the estimated hazard ratios from the basic Cox PH model. Of the nine covariates included, two are statistically significant at the 1\% level and one at the 5\% level. The remaining six covariates, including age, experience, team controls, the Giro d'Italia dummy, and the year trend, are not statistically distinguishable from zero.

\begin{table}[H]
\centering
\caption{Basic Cox proportional hazard model. The outcome is rider abandonment. Hazard ratios above 1 indicate higher abandonment risk; below 1 indicate lower risk. Significance: *** $p<0.001$, ** $p<0.01$.}
\label{tab:cox_basic}
\small
\renewcommand{\arraystretch}{1.3}
\begin{tabular}{@{} lrrrr @{}}
\toprule
Variable & Hazard Ratio & 95\% CI & $p$-value & \\
\midrule
Age                 & 0.998 & [0.988, 1.009] & 0.736 & \\
Experience          & 1.004 & [0.994, 1.015] & 0.424 & \\
Previous DNF rate   & 1.549 & [1.316, 1.822] & $<$0.001 & *** \\
GC contender        & 1.156 & [1.040, 1.286] & 0.008 & ** \\
Team size           & 0.992 & [0.903, 1.090] & 0.872 & \\
Team average age    & 0.987 & [0.963, 1.011] & 0.291 & \\
Tour de France      & 0.865 & [0.790, 0.946] & 0.002 & ** \\
Giro d'Italia       & 1.065 & [0.977, 1.161] & 0.150 & \\
Year (centred)      & 0.999 & [0.989, 1.009] & 0.836 & \\
\bottomrule
\multicolumn{5}{l}{\footnotesize Baseline race: Vuelta a Espa\~{n}a. $N = 9{,}045$ rider-race observations.}
\end{tabular}
\end{table}

The strongest predictor is the previous DNF rate (HR $= 1.549$, $p < 0.001$), implying that a rider who has always abandoned faces roughly 55\% higher hazard than one who has never abandoned.\footnote{The coefficient is reported per unit (0\% to 100\% historical DNF rate). For a one-SD increase ($\approx$22.6 pp), the effect is approximately $1.549^{0.226} \approx 1.10$.} This is consistent with two complementary interpretations: past abandonment reveals a ``fragile type'' with higher expected continuation costs, and from a signalling perspective \parencite{kahn2000labormarket}, a high DNF rate is a negative performance signal that may reflect career-stage dynamics.

The GC contender dummy (HR $= 1.156$, $p = 0.008$) indicates approximately 16\% higher abandonment hazard for GC contenders. Although seemingly counterintuitive, this is reconcilable with the incentive framework: GC contenders ride near their physiological ceiling and are more vulnerable to accumulated fatigue, and teams may strategically withdraw GC contenders who have fallen out of contention to preserve them for future races. In the tournament theory framework of \textcite{lazear1981rank}, the expected prize diminishes sharply once a contender falls too far behind the leader, making exit the rational response.

The Tour de France dummy (HR $= 0.865$, $p = 0.002$) implies roughly 14\% lower abandonment hazard than the Vuelta. This is consistent with the event-value hypothesis: the Tour offers higher prize money, greater media coverage, and more UCI ranking points, raising the economic returns from continued participation. This connects to the media economics framework in the course material (Lecture F3): broadcasting revenues and sponsor contracts are tied to participation and visibility \parencite{blair2012sports}. The Giro d'Italia coefficient is positive but not significant. The insignificance of age, experience, and team controls suggests these do not independently predict abandonment once DNF history and GC status are accounted for. Figure~\ref{fig:hazard_ratios} in the Appendix displays the hazard ratios in a forest plot.

\subsection{Time-Varying Cox Model}

Table~\ref{tab:cox_timevarying} extends the basic model by adding stage-level covariates that vary within a race.

\begin{table}[H]
\centering
\caption{Time-varying Cox proportional hazard model. Stage-type dummies and standardised distance are entered as time-varying covariates. Significance: * $p<0.05$, $\dagger$ $p=0.05$.}
\label{tab:cox_timevarying}
\small
\renewcommand{\arraystretch}{1.3}
\begin{tabular}{@{} lrrrr @{}}
\toprule
Variable & Hazard Ratio & 95\% CI & $p$-value & \\
\midrule
Age                 & 1.000 & [0.996, 1.003] & 0.923 & \\
Experience          & 1.000 & [0.997, 1.004] & 0.814 & \\
Previous DNF rate   & 1.067 & [1.000, 1.137] & 0.050 & $\dagger$ \\
GC contender        & 1.020 & [0.980, 1.061] & 0.341 & \\
Team size           & 1.000 & [0.972, 1.029] & 0.995 & \\
Tour de France      & 0.981 & [0.951, 1.011] & 0.215 & \\
Giro d'Italia       & 1.015 & [0.985, 1.047] & 0.335 & \\
Year (centred)      & 1.000 & [0.997, 1.003] & 0.943 & \\
Stage: hilly        & 0.997 & [0.964, 1.032] & 0.885 & \\
Stage: mountain     & 1.034 & [1.004, 1.064] & 0.025 & * \\
Stage: indiv.\ TT   & 0.965 & [0.915, 1.019] & 0.198 & \\
Stage: team TT      & 0.965 & [0.879, 1.059] & 0.454 & \\
Distance (z)        & 1.012 & [0.997, 1.027] & 0.111 & \\
\bottomrule
\multicolumn{5}{l}{\footnotesize Baseline stage type: flat. Baseline race: Vuelta a Espa\~{n}a.}
\end{tabular}
\end{table}

Mountain stages are the only stage type that significantly predicts abandonment (HR $= 1.034$, $p = 0.025$), consistent with a marginal-cost-of-effort interpretation: mountain stages impose higher physical demands and raise the instantaneous probability of abandonment on that day.

Once stage-type controls are included, several rider-level predictors lose significance. The GC contender dummy falls from HR $= 1.156$ ($p = 0.008$) to HR $= 1.020$ ($p = 0.341$), and the Tour de France dummy from HR $= 0.865$ ($p = 0.002$) to HR $= 0.981$ ($p = 0.215$). This suggests that part of what the basic model attributed to race identity and rider type was actually the correlation between these variables and stage-level difficulty. The previous DNF rate shrinks to HR $= 1.067$ (marginally significant at $p = 0.050$), partly because the time-varying model estimates effects on the daily rather than race-level hazard.

Taken together, the two models tell a consistent story. Riders with higher historical DNF rates face elevated abandonment hazard, consistent with higher expected continuation costs. The Tour de France generates lower abandonment, consistent with higher incentives from event value. Mountain stages raise the marginal cost of continuation on a given day. Because no direct financial variables are observed, these results should be read as suggestive evidence that economic incentives and costs shape abandonment patterns, not as a definitive causal decomposition.
